\documentclass[11pt]{letter}

\usepackage{graphicx}
\usepackage{amsmath,amssymb}
\usepackage{xspace}

% adjust margins
\usepackage{geometry}
\geometry{margin={2cm,2.5cm}}

% Self defined commands
\newcommand{\degC}{$^{\circ}$C\xspace}
\newcommand{\dC}{$\delta^{13}$C\xspace}
\newcommand{\dO}{$\delta^{18}$O\xspace}
\newcommand{\SrSr}{$^{87}$Sr/$^{86}$Sr\xspace}
\newcommand{\permil}{\textperthousand\xspace}
\newcommand{\UPb}{$^{206}$Pb/$^{238}$U\xspace}

\begin{document}

% add logo and address
\begin{minipage}{0.35\textwidth}
\includegraphics[width=\linewidth]{Figures/UCBerkeley_wordmark_blue.eps}
\end{minipage}
\hfill
\begin{minipage}{0.6\textwidth}\raggedleft
Yuem Park\\
Ph.D. candidate\\
307 McCone Hall\\
University of California, Berkeley\\
Berkeley, CA 94720\\
yuempark@berkeley.edu
\end{minipage}

\vspace{0.5cm}
\hrulefill
\vspace{-0.1cm}

\null\hfill \today

Dear GSA Bulletin Science Editors,

This letter accompanies our submission of a manuscript for publication in GSA Bulletin entitled `The lead-up to the Sturtian Snowball Earth: Neoproterozoic chemostratigraphy time-calibrated by the Tambien Group of Ethiopia.'

The first of the severe Cryogenian `Snowball Earth' glaciations - the ca. 717-660~Ma Sturtian Glaciation - terminates an extended interval ($\sim$1.5 billion years) of ice-free conditions, and thus represents a fundamental shift in Earth's climatic state. However, the mechanisms responsible for initiating such a dramatic shift remain enigmatic. One important tool that can be used to constrain surface conditions and better understand the mechanisms behind the Sturtian Glaciation are chemostratigraphic data recorded in marine carbonates. For example, isotope proxies for the global carbon cycle (\dC) and weathering fluxes (\SrSr) suggest major perturbations in these systems leading into the Sturtian Glaciation. However, due to a lack of precise constraints on the timing, duration, and rates of these perturbations, it is unclear if/how they relate to the initiation of severe glaciation. Furthermore, without direct temporal constraints on chemostratigraphic datasets from sedimentary sequences in separate basins around the world, the interpretation that these perturbations to the surface environment are global, rather than local, cannot be tested.

The Tonian-Cryogenian Tambien Group of northern Ethiopia is a mixed carbonate-siliciclastic sequence of sedimentary rocks that culminates in diamictite associated with the Sturtian Glaciation. The presence of intercalated tuffs suitable for high precision U-Pb CA-ID-TIMS geochronology, as well as the completeness of the stratigraphy leading into the glaciation, makes these sedimentary rocks a unique and ideal target for temporally constraining physical and isotopic stratigraphic data sets leading up to the Sturtian Glaciation. In this manuscript, we utilize the recent discovery of vast exposures of Sturtian diamictite and underlying Tambien Group strata in the Samre Fold-Thrust Belt to develop new lithostratigraphic, geochemical, and geochronologic data from this interval. U-Pb dates from tuffs just below the diamictite demonstrate that the Samre Fold-Thrust Belt preserves pre-glacial carbonate stratigraphy that is far closer to the onset of the Sturtian Glaciation than has been identified anywhere else in the world, and therefore presents a unique opportunity to geochemically constrain surface processes that were occurring at this time.

Combining our \dC, \SrSr, and U-Pb data with that from several other basins around the world, we have constructed the most temporally well-constrained composite Tonian carbon and strontium isotope curves to date. The composite \dC dataset suggests that the sharpest perturbations to the carbon cycle are globally synchronous and are not associated with onset of the Sturtian Glaciation. The composite \SrSr dataset suggests that enhanced subaerial weathering of juvenile lithologies, and thus enhanced consumption of CO$_{2}$, began $\sim$50~My prior to the initiation of the Sturtian Glaciation. We also analyze the timing and paleolatitude of large igneous province eruptions and arc accretion events relative to the \SrSr curve, and argue that the accretion of Arabian-Nubian Shield volcanic arcs are an overlooked but important contributor to the long-term changes in tectonic boundary conditions that set the stage for the initiation of the first Neoproterozoic Snowball Earth.

We believe that this contribution represents a significant advance in the study of Neoproterozoic Earth history and will be of wide interest to the readership of GSA Bulletin.

Sincerely, and on behalf of my co-authors,

\vspace{0.2cm}
\includegraphics[width=0.25\linewidth]{Figures/signature.pdf}
\vspace{0.2cm}

Yuem Park

\vspace{1cm}

Suggested reviewers:

\begin{itemize}
\item Carol Dehler (carol.dehler@usu.edu)
\item Galen Halverson (galen.halverson@mcgill.ca)
\item Graham Shields (g.shields@ucl.ac.uk)
\item Emmy Smith (efsmith@jhu.edu)
\item Michael Beyth (mbeyth@gsi.gov.il)
\item Charles Verdel (c.verdel@uq.edu.au)
\item Francis Macdonald (fmacdon@eps.harvard.edu)
\item Melina Macouin (melina.macouin@get.omp.eu)
\item Alan Rooney (alan.rooney@yale.edu)
\item David Jones (djones@amherst.edu)
\end{itemize}

\vspace{1cm}

\footnotesize
Note for the Managing Editor: We ask that the submitted PDF of the main text with integrated figures be utilized for review. The manuscript with references was drafted using LaTEX+BibTEX and we wish to keep it in that format through revisions. While we wish GSA accepted this format, we can convert the manuscript into .doc format following revisions of the manuscript.

\end{document}